\section{Continuing group search across word representations}
\label{sec:adaptive-group}

\subsection{Why the representation should depend on progress}

The preceding measurements show two distinct regimes. Explicit relators
are faster on the current finite knot corpus, while compressed words can
represent intermediate strings beyond an explicit-letter allowance.
The new optional adaptive stage starts with the explicit producer and
changes representation once when the next substitution would exceed a
chosen threshold. It continues the existing search state, preserving the
certificate prefix and remaining resources. This implements a form of
introspective search without requiring compressed arithmetic on every
small initial relation.

Let $C$ be the hard explicit-letter allowance and $T\le C$ the switching
threshold. By default $T=C$. A caller may request an earlier nonnegative
threshold; values above $C$ are clamped to $C$. Initial presentation
construction must still satisfy the hard allowance, including the
$4n$-letter bound before crossing relations are simplified. Adaptivity
does not bypass that check.

\subsection{Predicting an allocation before changing state}

Suppose the selected defining relator $w$ has length $\ell$ and contains
the eliminated generator $g$ exactly once. Let $L$ be the total relator
length and $c_g$ the total number of occurrences of $g^{\pm1}$. The
replacement for $g$ has length $\ell-1$. Removing $w$ and substituting
into the remaining words would allocate
\[
 E=L-\ell+(c_g-1)(\ell-2)
\]
letters before free reduction. This is an exact raw substitution length,
including the case $\ell=1$, where all remaining occurrences of $g$ are
deleted. If $E>T$, the producer switches before deleting $w$, constructing
the replacement, changing the generator set or recording the pending move.
Otherwise the existing explicit substitution and reduction proceed.

A Whitehead image has length at most three, so $3L$ is a conservative
allocation bound. A selected shortening Whitehead move triggers conversion
if $3L>T$, again before applying or recording it. Its final reduced length
may be much smaller: the trigger controls the raw allocation allowance,
not a theorem that the move necessarily expands the presentation. Relator
overlap moves strictly decrease total length and retain their existing
explicit implementation during this phase.

The producer converts each current relator to an SLP root while preserving
its slot, including empty slots. Generator names and the surviving set are
unchanged. It passes the same move list to the compressed continuation;
there is no second reconstruction of the initial presentation and no
replay of the producer's completed moves. The continuation may recompute
the pending search choice, which is charged normally. This duplication
can occur at only one handoff. Once compressed, the producer never returns
to explicit search, avoiding repeated conversion near the threshold.

\subsection{Soundness and the resource boundary}

Conversion changes only representation, so the group presented by the
current relators is unchanged. Every subsequent move uses the same Tietze,
Whitehead or relator-overlap justification as before. The resulting
certificate format is unchanged and remains bound to the original input
PD. Acceptance still uses the cyclic-knot-group criterion of
Section~\ref{sec:group-certificates}.

If no conversion occurred, independent replay remains explicit unless the
caller separately requested compressed verification. If conversion did
occur, compressed replay is selected automatically: literal replay might
exhaust the very allowance that prompted the producer to switch. In both
cases the checker independently reconstructs the full original Wirtinger
presentation and verifies every recorded move from the beginning. It does
not accept the handoff state as evidence. The two implementations still
share the underlying compressed-word kernel when compressed replay is
selected, as described in the previous section.

The explicit prefix consumes part of the original work allowance. The
SLP arena receives exactly the remainder, and conversion itself is charged
to it. The continuation does not obtain a fresh work budget or a new wall
deadline. Its node allowance is the usual compressed node/assertion cap.
Independent replay retains its separately specified work allowance, while
search and verification share the local wall deadline. A depleted work
allowance, node-limited conversion, local timeout or stalled continuation remains
inconclusive; global caller cancellation propagates.

Evidence reports whether conversion occurred and records its trigger,
the number of preceding moves, surviving rank, current and projected
lengths, threshold, work already spent and work remaining. These are
diagnostics, not additional certificate axioms. A regression test leaves
exactly one work unit at handoff and checks that conversion exhausts it
rather than silently refilling the budget.

\subsection{A bound on the explicit prelude}

Although no quasi-polynomial bound for the full recognizer is established,
the explicit portion has a useful conditional bound. Let $r$ be the initial
generator rank, let $L_0$ be its initial reduced total relator length, and
put $M=\max\{L_0,T\}$. There are at most $r-1$ generator eliminations.
Between eliminations, each accepted Whitehead or overlap move decreases
the nonnegative integer total length by at least one. An elimination
performed explicitly has raw output length at most $T$, and free reduction
cannot increase that length. Every interval between eliminations therefore
starts at length at most $M$ and contains at most $M$ shortening moves.
Including eliminations and the possible final switch decision gives
\[
 O\bigl(r(M+1)\bigr)
\]
explicit move iterations. Counts, suffix-automaton overlap searches,
minimum cuts and literal substitutions are polynomial in the current
explicit input size and rank. Conversion is polynomial in the current
length too, and occurs once. Thus a polynomially bounded threshold gives
a polynomially bounded explicit prelude on inputs admitted by its initial
resource checks.

This does not establish a polynomial or quasi-polynomial recognizer. The
compressed continuation still needs a completeness argument and bounds on
both grammar growth and the number of moves for the reachable knot
presentations. Moreover, finding
positive cyclic-group certificates quickly would not itself supply a fast
conclusive route on nontrivial knots. The surrounding recognizer retains
its previous complete exponential fallback. The policy preserves progress
and limits representation overhead; it does not manufacture a stronger
worst-case guarantee from its current backends.

\subsection{Interface and validation}

The CLI flag \code{--group-adaptive} enables this group stage with $T=C$.
\code{--group-switch-letters T} also enables it and selects an earlier
threshold. Python recognition accepts:
\begin{verbatim}
recognize(diagram, use_group=True, group_adaptive=True,
          group_switch_letters=T)
\end{verbatim}
For standalone group calls, enable \code{adaptive\_search}.
Always-compressed search and
adaptive search are mutually exclusive. Replay-only compression may be
combined with either policy. The default recognition pipeline is unchanged.

Tests show that the default adaptive policy preserves all fifteen corpus
certificates exactly and never constructs a search arena on those paths.
Earlier-switch traces pass both independent replayers. Tests force both
elimination and Whitehead handoffs, verify the preserved trace prefix,
prohibit reconstruction at conversion, check resource and cancellation
behavior, reject invalid option combinations and exercise CLI routing.
Known nontrivial knots remain inconclusive at this positive-certificate
stage. These finite checks support the invariants above without implying
completeness of the group heuristic.

The full integrated suite passes all 653 tests with the optional Regina
dependency installed in 126.251 seconds. Its output is preserved in
\path{data/adaptive-group-integrated-tests.txt}.

\subsection{Whole-query and matched-cap measurements}

The 8 October 2026 audit runs five interleaved whole-query arms on the same
fifteen real diagrams: explicit search, an identical explicit control,
adaptive search with $T=C=200{,}000$, adaptive search with $T=4{,}096$ and
$C=200{,}000$, and always-compressed search. Every arm includes fresh PD
validation, all recognition stages and independent group-certificate replay.
There are five measured rounds after a warm-up, giving 375 whole queries.
Search and replay each have twenty million work units; the local group
allowance is fifteen seconds, the global allowance eighteen seconds,
and the node allowance 100,000. Local tests do not run concurrently.

All queries finish with UNKNOT and group evidence. The default adaptive
policy does not switch on any of these inputs. The earlier threshold
switches only on Gordian; its fourteen small-case paths stay explicit.
The small-case column is the sum of the fourteen per-input medians:
\begin{center}
\begin{tabular}{lrr}
\toprule
Search policy & Small cases (ms) & Gordian (s)\\
\midrule
Explicit & 52.54 & 1.291\\
Explicit control & 52.18 & 1.251\\
Adaptive, $T=C$ & 52.00 & 1.242\\
Adaptive, $T=4{,}096$ & 52.05 & 1.641\\
Always compressed & 190.37 & 1.970\\
\bottomrule
\end{tabular}
\end{center}
The no-switch adaptive results are near the explicit controls; these small
differences are not presented as an algorithmic speedup. Early conversion
costs more on Gordian than letting its explicit path finish. Accordingly,
the default trigger remains the hard letter ceiling, not the earlier
threshold tested here. The optional policy is primarily a capacity
extension that retains the faster representation while it fits.

A second experiment lowers the hard letter cap on four real-input cases.
It measures fresh PD validation and the group stage alone, including
independent replay, without the surrounding recognition filters or fallback.
The four arms are explicit, its control, adaptive with $T=C$, and
always compressed. Five measured rounds after a warm-up give 80 probes.
The explicit arms exhaust the letter allowance in all forty measured
probes. Both adaptive and always-compressed search produce independently
verified certificates in all forty of their probes:
\begin{center}
\begin{tabular}{lrrr}
\toprule
Input & $C$ & Adaptive median (s) & Compressed median (s)\\
\midrule
Survivor 00 & 64 & 0.01093 & 0.01380\\
Survivor 02 & 60 & 0.01102 & 0.01285\\
Gordian & 4,096 & 3.72104 & 3.41667\\
Gordian & 8,192 & 1.15899 & 1.23270\\
\bottomrule
\end{tabular}
\end{center}
The explicit exhaustion times are recorded in the raw artifact, but are
not completed recognition times and are not used in speedup ratios. This
experiment establishes additional certificate capacity under matched caps,
not newly decidable knots: the overall recognizer already has its complete
fallback, and sufficiently generous explicit limits certify these examples.

Hard and soft thresholds have different effects. With $T=4{,}096$ but
$C=200{,}000$, Gordian switches after 127 completed moves from 3,228 current
letters before an 8,098-letter raw substitution. The continuation still
has the large overlap allowance and finishes with 141 total moves and
7,270 allocated nodes. When the hard cap is also 4,096, the switch point
is the same, but the compressed continuation skips 75 oversized overlap
probes. It follows a longer path of 220 moves and allocates 76,758 nodes.
At hard cap 8,192, conversion occurs after 128 moves; the continuation
skips five overlap probes and finishes with 149 moves and 21,279 nodes.
The cap therefore changes more than memory use: it can change the chosen
sequence of valid transformations and its cost. No threshold is claimed
uniformly optimal.

The reproducible driver is \path{fast/benchmark_adaptive_group.py} and the
raw artifact is \path{fast/results/adaptive_group_20261008.json}.
The evidence supports a correct one-way continuation mechanism and measured
capacity gains. It leaves the general sub-exponential recognition question
open.
